\section{Introduction}
\label{sec:introduction}

 India's Business Responsibility and Sustainability Reporting (BRSR) framework gives environmental, social, and governance (ESG) disclosures a common structure. The Securities and Exchange Board of India (SEBI) made BRSR reporting mandatory in 2021 for the top 1,000 listed companies by market capitalisation, and the rule applies from financial year 2022--23. The framework rests on nine principles of responsible business conduct. Essential indicators must be reported, but leadership indicators are voluntary \cite{sebi2021}. Every year, each company files one complete document, called a \emph{report}. A common format makes reports easier to organise. The format alone, however, cannot show whether every answer is complete or correct.

Reviewing one report by hand is a slow job with several steps. The analyst first finds a \emph{text segment}, which is the selected portion of the report's text that is analysed, and places the text segment under a topic. Then the analyst looks for a related number. This number must be checked for the reporting year, the unit, and the part of the business that the number covers. Comparing companies adds more difficulty. Business activities differ, so different reporting items apply to each company. Across hundreds of reports, these steps take a great deal of time. Take the claim \emph{water use decreased} as an example. The claim can be assessed only after a matching water measure, the business boundary, and the comparison period are found.

Automated text analysis can reduce this workload. ClimateBERT finds climate-related language and some disclosure categories. ChatReport retrieves text from reports and answers reporting questions, with a link to the source for each answer. ESGReveal pulls structured information out of narrative reports \cite{webersinke2021,ni2023,zou2025}. Methods like these make large collections of reports easier to organise. Even so, three questions remain once a connected workflow is built for structured Indian reports.

Three questions are still open. The first is \emph{how to connect a claim with evidence}. A phrase detector can find a promise, and a text-similarity model can tell that a sentence is about the environment. Neither tool checks the number behind the sentence. That number must fit the meaning of the claim, the reporting period, and the reporting context before the number counts as support.

The next question is \emph{how to organise training labels}, meaning the correct answers from which a model learns. Labelling every text segment by hand would be slow. Tags and concept maps in the report offer a repeatable alternative, and the source of each label is recorded for every text segment.

The last question is \emph{how to interpret a company score}. A score means little unless the reader knows how the score was produced. This requires stating the calculation, the items included, the sector group, and the reference source. GRASP brings these requirements together in one connected BRSR workflow with clearly defined tests \cite{calamai2025,garg2025}.

Each of these questions gets a direct answer in GRASP. For the link between claim and evidence, text segments and structured facts sit under the same report. A \emph{fact} is one reported value together with the concept, unit, and period. When a language check flags a text segment, the analyst can open the text segment directly. The numerical and assurance checks work the same way, because each one keeps the linked sentence and the reported item ready for inspection. For the labels, three label sources are recorded, and results are compared by source, by company, and by reporting question left out of training. For the score, four components are shown, and each company is compared with others in the same recorded sector. Together, these steps give GRASP links to filed evidence, training labels taken from the structure of reports, a detailed evaluation, and a score with a visible breakdown.

The study provides a dataset of 42,409 text segments, each stored as a dataset row with a documented weak label. A \emph{weak label} is a principle label assigned automatically from the structure of the report, instead of by a person reading the text. Separate companies are used for training, validation, and testing. The classifier is assessed by label source and on concepts left out of training. The language detector is evaluated with fixed settings, and the numerical checks supply linked review records. Company-score results are presented at four levels, from one worked example to the full financial-sector group and then all recorded sectors. For every model result, human-score comparison, and service check, the evaluation states both the data used and the reference for comparison.

The rest of the paper moves from earlier research to practical use. Section~\ref{sec:literature} reviews related methods and Indian disclosure studies. The proposed solution comes in Section~\ref{sec:proposed}, which traces the path from report data to review outputs. Section~\ref{sec:methodology} then describes the seven modules and the score calculations in detail. After that, Section~\ref{sec:evaluation} sets out the test data and the audited model results. Section~\ref{sec:results} shows the saved outputs at four levels, from one company to all sector groups. The paper ends in Section~\ref{sec:conclusion} with the main findings and future work.